\section{Related work}
\label{sec:related}
% ~0.75 page. Each sentence maps to a LIT-xx claim in report/claims.json.
% Primary arXiv identities were checked in a later audit; consult docs/CITATION_AUDIT.md.

% LIT-01
Dynamical decoupling of open quantum systems was introduced by \citet{Viola1999}.
% LIT-02
\citet{Uhrig2007} proposed optimized $\pi$-pulse sequences for keeping a qubit coherent,
% LIT-03
and \citet{Cywinski2008} analysed how to enhance dephasing times in superconducting qubits.
% LIT-04
\citet{Biercuk2009} studied optimized DD in a model quantum memory.
% LIT-05, LIT-06
Biercuk, Doherty, and Uys described DD optimization as a filter-design
problem \citep{BiercukDohertyUys2011}, while Green and coauthors developed
higher-order filter methods for driven gates \citep{GreenUysBiercuk2012}.
These approaches directly motivate deterministic model-based objectives,
which differ from the finite-trajectory CMA-ES objective used here.
DD has also been used for noise spectroscopy: \citet{Bylander2011} with a superconducting flux
qubit and \citet{Alvarez2011} for colored noise.
% LIT-07
Robust DD was studied by \citet{Souza2011},
% LIT-08
and \citet{Ezzell2023} surveyed DD performance on superconducting qubits.
% LIT-09, LIT-10, LIT-11
Empirical or hardware-tailored learning approaches include \citet{Rahman2024}, \citet{Tong2025}
and \citet{Tong2026}.

\paragraph{Empirical versus model-based}
% LIT-10 (title only)
\citet{Rahman2024}, \citet{Tong2025}, and \citet{Tong2026} study
empirical or hardware-tailored dynamical decoupling. In particular,
\citet{Rahman2024} optimizes rotational gates for superconducting
hardware, while \citet{Tong2026} studies empirical suppression in dynamic
circuits. The present approach instead uses a fitted or assumed synthetic
noise model as the optimization objective. Such a model-based approach
reduces direct hardware calls but depends on model fidelity.

\paragraph{Positioning}
% No LIT claim: describes this project only. Novelty claim deliberately limited.
This project is a model-based, synthetic-only study: sequences are optimized and evaluated
against a simulated noise model, and no quantum hardware has been used.
Its contribution is limited to a pre-registered, held-out synthetic benchmark of model-based DD
optimization with explicit noise floors.
Its main finding is negative: under the pre-registered synthetic model, CMA-ES did not
provide a practically meaningful generalizable improvement over XY4/XY8 on held-out noise (Sec.~\ref{sec:results}).
A synthetic benchmark of this kind does not test whether its conclusions transfer to hardware.
